\section{The collar condition is computable}\label{sec:collar}

\subsection{Weak feature size}
Let $M\subset\R^3$ be compact, $d(p)=\dist(p,M)$, and let $\Gamma(p)$ be the set of nearest points of $p$ in $M$.
A point $p\notin M$ is \emph{critical} if $0\in\conv\{p-x:x\in\Gamma(p)\}$. The \emph{weak feature size}
$\wfs(M)$ is the infimum of the positive critical values of $d$~\cite{ChazalLieutier07}. Unlike the reach, it
stays positive at sharp edges and small fillets. A surface with a concave edge has reach $0$, and a fillet of
radius $r$ caps the reach at $r$.

\begin{lemma}[Collar]\label{lem:collar}
Let $M$ be a closed PL surface in $\R^3$, and let $0<r<\wfs(M)$.
\begin{enumerate}
\item[(a)] The closed $r$-neighbourhood $\bar N_r(M)$ is PL-homeomorphic to $M\times[-1,1]$, with $M$ going to
$M\times\{0\}$.
\item[(b)] Suppose every zone of every cell of a shell lies in $N_d(M)$ for some $d<r$, $M$ is the image of a
carried PL embedding of that shell, and the neighbourhoods $\bar N_r(M_i)$ of different shells are pairwise
disjoint. Then the collar condition of Section~\ref{sec:uniqueness} holds with $U_i=\bar N_r(M_i)$.
\end{enumerate}
\end{lemma}
\begin{proof}[Proof sketch]
\begin{itemize}
\item \emph{Small radii.} For small $\delta>0$, $\bar N_\delta(M)$ is a regular neighbourhood of $M$ relative to
$M$, by Durfee's theorem on neighbourhoods of semialgebraic sets~\cite[Props.~3.5,~3.7]{Durfee83}. A regular
neighbourhood of a closed orientable surface in $\R^3$ is a product.
\item \emph{Larger radii.} The region $\{\delta\le d\le r\}$ contains no critical point, so it is a product of
a level set with $[\delta,r]$, by the isotopy lemma of Chazal and Lieutier~\cite[Prop.~3.4]{ChazalLieutier07}.
\item \emph{Gluing.} Gluing the two products gives (a). Part (b) follows directly.
\end{itemize}
\draftnote{The application of~\cite[Prop.~3.4]{ChazalLieutier07}, which is stated for distance to the
complement of a bounded open set, to $d=\dist(\cdot,M)$ is spelled out in the repository notes (review of Lemma~1); it should be written out here. The PL (rather than
topological) form of (a) also needs a sentence.}
\end{proof}

\subsection{An explicit lower bound for polyhedral surfaces}
For $a\in M$, the \emph{conical radius} $\rho(a)$ is the supremum of the $r>0$ such that $M\cap B(a,r)=C\cap
B(a,r)$ for some closed cone $C$ with apex $a$.

\begin{lemma}[Cone lemma]\label{lem:cone}
Let $a\in M$, $r<\rho(a)$, and $p\notin M$ with $|p-a|<r/2$. Then every nearest point $x\in\Gamma(p)$ satisfies
$\langle p-x,\,p-a\rangle=|p-x|^2$. In particular $p$ is not critical.
\end{lemma}
\begin{proof}
\begin{itemize}
\item \emph{Nearest points lie in the ball.} $d(p)\le|p-a|<r/2$, so every $x\in\Gamma(p)$ has $|x-a|<r$.
\item \emph{The ray through $x$ stays in $M$.} The points $y_s=a+s(x-a)$, for $s$ near $1$, lie in the cone and
in $B(a,r)$, hence in $M$. So $s\mapsto|p-y_s|^2$ has a local minimum at $s=1$, which gives $\langle
p-x,x-a\rangle=0$.
\item \emph{Conclusion.} $\langle p-x,p-a\rangle=|p-x|^2+\langle p-x,x-a\rangle=d(p)^2>0$ for all
$x\in\Gamma(p)$. So $0$ is not in the convex hull of the vectors $p-x$.
\end{itemize}
\end{proof}

\begin{corollary}\label{cor:lambda}
Let $\lambda(M)=\inf_{x\in M}\sup_{a\in M}\bigl(\rho(a)/2-|x-a|\bigr)$. Then $d$ has no critical point $p$ with
$0<d(p)<\lambda(M)$, so $\wfs(M)\ge\lambda(M)$. If $M$ is a finite union of triangles, $\lambda(M)>0$.
\end{corollary}
\begin{proof}
\begin{itemize}
\item \emph{The bound.} Given $p$ with $d(p)<\lambda$ and $x\in\Gamma(p)$, pick $a$ with
$\rho(a)/2-|x-a|>d(p)$. Then $|p-a|<\rho(a)/2$, and Lemma~\ref{lem:cone} applies.
\item \emph{Positivity.} For a polyhedron, $\rho(a)>0$ for every $a$. The function $g(x)=\sup_a(\rho(a)/2-|x-a|)$
is $1$-Lipschitz and finite, since $\rho\le\operatorname{diam}M$, and $g(x)\ge\rho(x)/2>0$. $M$ is compact, so
$\min g>0$.
\end{itemize}
\end{proof}

Lemma~\ref{lem:cone} and the pointwise form of Corollary~\ref{cor:lambda} are machine-checked in Lean~4 with
Mathlib, for subsets of an arbitrary real inner product space (Appendix~\ref{sec:lean}).

\subsection{A certified lower bound for $\wfs$}
For a triangulated reference surface we certify ``no critical value in $(0,L)$'' with a best-first octree:
\begin{itemize}
\item \emph{Boxes near $M$} are certified by Lemma~\ref{lem:cone}, using a lower bound for $\rho(a)$ at a
candidate point $a$. The bound comes from the distance from $a$ to the nearest feature (planar region, boundary
segment, vertex) that does not contain $a$.
\item \emph{Boxes farther away} are certified by a direction test: a small linear program finds a unit vector $u$
with $\langle u,p-x\rangle>0$ for all candidate nearest points $x$ of all points $p$ of the box. The candidates
are the projections of the box centre onto the relevant features.
\item \emph{Stopping rule.} The search stops at the first box that is too small to split further and cannot be
certified. Its distance from $M$ is the certified bound $L$.
\end{itemize}
The computation is in double precision, not interval arithmetic. Table~\ref{tab:wfs} lists the results.

\begin{table}[h]
\centering
\begin{tabular}{lcc}
\toprule
Part & exact $\wfs$ & certified lower bound \\
\midrule
cube, side 10 & 5 & 4.530 \\
plate $4\times4\times1$ & 0.5 & 0.436 \\
plate $4\times4\times0.3$ & 0.15 & 0.137 \\
L-shape (wall 2) & 1 & 0.894 \\
L-shape, inner fillet 0.2 & 1 & 0.894 \\
block with slot of width 0.2 & 0.1 & 0.086 \\
cylinder $R=3$, $H=10$ & 3 & 2.512 \\
\bottomrule
\end{tabular}
\caption{Certified lower bounds for $\wfs$ of the triangulated parts. The inner fillet caps the reach at
$0.2$ but leaves the certified $\wfs$ bound unchanged.}\label{tab:wfs}
\end{table}
